\chapter{VGC Library}

The VGC library is the assembly programmer's thin layer over Nova's Virtual
Graphics Controller. It does not hide the hardware model. It gives names to the
registers, exposes the command protocol, and provides small helper routines for
the operations that almost every graphics program needs.

\begin{novanote}
Use \novacode{vgc.inc} for declarations and constants. Include
\novacode{vgc.s} when the program wants the helper routines, or include
\novacode{vgc\_wait.s} when an overlay only needs the command wait primitive.
\end{novanote}

\begin{lstlisting}
.include "vgc.inc"

      ; program code

.include "vgc.s"
\end{lstlisting}

\section{What The VGC Is}

The VGC is Nova's display custom chip. From the 6502 side it is a memory mapped
device at \novacode{\$A000-\$A0FF}. Inside the device are separate engines for
video timing, text, bitmap graphics, sprites, raster-time register writes, IRQs,
and off-bus video memory.

\begin{center}
\novasvg[width=\textwidth]{build/diagrams/vgc-architecture}
\end{center}

The important mental model is that the CPU does not draw pixels directly into a
linear framebuffer in main RAM. Instead, code talks to the VGC through three
paths:

\begin{itemize}
\item Direct registers for state such as mode, color, cursor position, frame
      counter, sprite attributes, palette mode, and IRQ latches.
\item The command port at \novacode{\$A010-\$A01F} for drawing primitives,
      sprite shape commands, copper list operations, and byte-level memory
      commands.
\item The VGC memory port, DMA, blitter, FIO, and pager paths for bulk access to
      character RAM, color RAM, the graphics bitmap, sprite shape RAM, and text
      attribute RAM.
\end{itemize}

\section{Components}

\begin{description}
\item[VESTA timing] owns the scan position, frame counter, vblank boundary, and
video sync. The \novacode{vgc\_vsync} helper waits for \novacode{VGC\_FRAME} at
\novacode{\$A008} to change.
\item[SCRIBE text] renders the 80 by 50 character layer from character, color,
font, and text-attribute memory. \novacode{VGC\_CHAROUT} accepts printable text
and control bytes such as form feed, carriage return, line feed, home, and
backspace.
\item[CANVAS graphics] stores the 320 by 200, 16-color bitmap. Drawing commands
such as \novacode{vgc\_line}, \novacode{vgc\_rect}, and \novacode{vgc\_paint}
operate on this memory through the ARTIST engine.
\item[ARTIST drawing] executes graphics commands from the command port. It uses
the current graphics color set by \novacode{VGC.GCOLOR} and clips drawing to the
bitmap.
\item[PIXIE sprites] renders 16 hardware sprites. Each sprite has an 8-byte
attribute block at \novacode{\$A040-\$A0BF}; shape pixels live in off-bus sprite
shape RAM.
\item[CONDUCTOR copper] stores raster events and writes selected VGC or sprite
registers when scanout reaches the programmed position.
\item[IRQ block] exposes enable, pending, force, and valid-source registers at
\novacode{\$A0F0-\$A0F3}. The VGC can raise IRQs for vblank, copper events, and
sprite collision sources.
\end{description}

\section{Command Helpers}

The command protocol is always the same: load \novacode{VGC\_P0} through
\novacode{VGC\_P14}, then write a command byte to \novacode{VGC\_CMD}. The
helper routines in \novacode{vgc.s} save you from remembering the command byte.
They do not all wait for completion.

\begin{center}
\novasvg[width=\textwidth]{build/diagrams/vgc-command-flow}
\end{center}

\begin{center}
\small
\begin{tabularx}{\textwidth}{@{}l>{\raggedright\arraybackslash}X@{}}
\toprule
\textbf{Helper pattern} & \textbf{Use} \\
\midrule
\novacode{vgc\_cmd} & Write A to \novacode{VGC\_CMD} and return immediately. \\
\novacode{vgc\_exec} & Write A to \novacode{VGC\_CMD}, then call
\novacode{vgc\_wait\_cmd}. \\
\novacode{vgc\_wait\_cmd} & Poll bit 0 of \novacode{VGC\_CMD} until the active
command completes. \\
\novacode{vgc\_mem\_read}, \novacode{vgc\_mem\_write} & Issue the memory
command and wait, because the caller needs the byte result or address update. \\
\novacode{vgc\_gcls}, \novacode{vgc\_line}, \novacode{vgc\_fill}, and other
primitive helpers & Issue the drawing command and return. Wait explicitly when
the next step depends on the result or reuses command parameters. \\
\bottomrule
\end{tabularx}
\end{center}

This distinction matters most on FPGA builds, where a clear, flood fill, sprite
clear, sprite copy, or text clear can span many pixel clocks. Avalonia's hosted
VGC often completes commands immediately, but portable code should synchronize
when it depends on command completion.

\section{Drawing From Assembly}

The following example selects graphics/sprite mode, clears the bitmap, sets a
draw color, and draws a diagonal line. It waits after the clear and after the
line because both operations modify the graphics plane.

\begin{lstlisting}
.include "vgc.inc"

draw_demo:
      LDX   #$03
      JSR   vgc_set_mode

      JSR   vgc_gcls
      JSR   vgc_wait_cmd

      LDA   #$07
      STA   VGC_P0
      JSR   vgc_gcolor

      ; LINE 0,0,319,199
      STZ   VGC_P0      ; x0 low
      STZ   VGC_P1      ; x0 high
      STZ   VGC_P2      ; y0 low
      STZ   VGC_P3      ; y0 high
      LDA   #$3F        ; x1 = $013F
      STA   VGC_P4
      LDA   #$01
      STA   VGC_P5
      LDA   #$C7        ; y1 = 199
      STA   VGC_P6
      STZ   VGC_P7
      JSR   vgc_line
      JSR   vgc_wait_cmd
      RTS

.include "vgc.s"
\end{lstlisting}

\section{Text And Display State}

The text helpers are direct register writes and return immediately:
\novacode{vgc\_cls}, \novacode{vgc\_locate}, \novacode{vgc\_set\_fg},
\novacode{vgc\_set\_bg}, \novacode{vgc\_set\_border},
\novacode{vgc\_set\_mode}, \novacode{vgc\_set\_font},
\novacode{vgc\_display\_on}, and \novacode{vgc\_display\_off}.

\novacode{vgc\_cls} writes form feed to \novacode{VGC\_CHAROUT}; it clears the
text plane and homes the cursor. \novacode{vgc\_locate} copies
\novacode{VGC\_P0} and \novacode{VGC\_P1} into the cursor X and Y registers.
Foreground, background, and border helpers take their color in X.

\begin{lstlisting}
      LDX   #15
      JSR   vgc_set_fg
      LDX   #0
      JSR   vgc_set_bg
      LDX   #11
      JSR   vgc_set_border

      LDA   #10
      STA   VGC_P0
      LDA   #6
      STA   VGC_P1
      JSR   vgc_locate
\end{lstlisting}

\section{Off-Bus VGC Memory}

Character RAM, color RAM, text attributes, graphics pixels, and sprite shapes
are VGC memory spaces. They are not mapped into the 6502 address range as normal
RAM. The library exposes \novacode{vgc\_mem\_read} and
\novacode{vgc\_mem\_write} for byte-level access, while larger transfers should
use the VRAM port, DMA, blitter, FIO graphics load/save, or pager routines.

\begin{center}
\novasvg[width=\textwidth]{build/diagrams/vgc-memory-paths}
\end{center}

\begin{center}
\small
\begin{tabularx}{\textwidth}{@{}cl>{\raggedright\arraybackslash}X@{}}
\toprule
\textbf{Space} & \textbf{Name} & \textbf{Use} \\
\midrule
1 & Character RAM & 4000 character bytes for the 80 by 50 text layer. \\
2 & Color RAM & 4000 packed background/foreground color bytes. \\
3 & Graphics bitmap & 64000 pixels, one byte per pixel, low nibble is color. \\
4 & Sprite shape RAM & 32768 bytes, 256 shapes at 128 bytes per shape. \\
7 & Text attribute RAM & 4000 style bytes; bit 0 is flash. \\
\bottomrule
\end{tabularx}
\end{center}

For isolated bytes, load \novacode{P0=space}, \novacode{P1/P2=offset},
\novacode{P3=data}, and \novacode{P4=flags}. Bit 0 of \novacode{P4}
auto-increments the offset after each read or write.

\begin{lstlisting}
      ; Write "OK" into character RAM at cell 0 using auto-increment.
      LDA   #VGC_PLANE_CHAR
      STA   VGC_P0
      STZ   VGC_P1
      STZ   VGC_P2
      LDA   #VGC_VRAM_AUTOINC
      STA   VGC_P4

      LDA   #'O'
      STA   VGC_P3
      JSR   vgc_mem_write

      LDA   #'K'
      STA   VGC_P3
      JSR   vgc_mem_write
\end{lstlisting}

For streams, prefer the VRAM data port. It uses the same space IDs and advances
the address in hardware when \novacode{VGC\_VCTRL} has
\novacode{VGC\_VRAM\_AUTOINC} set. Reads are latched: the first read from
\novacode{VGC\_VDATA} posts the selected byte and the second read returns it, so
read loops should prime the latch before consuming data.

\begin{lstlisting}
      LDA   #VGC_PLANE_GFX
      STA   VGC_VPLANE
      STZ   VGC_VADDRL
      STZ   VGC_VADDRH
      LDA   #VGC_VRAM_AUTOINC
      STA   VGC_VCTRL

      ; Stream pixels by writing VGC_VDATA.
      LDA   #$05
      STA   VGC_VDATA
      STA   VGC_VDATA
      STA   VGC_VDATA
\end{lstlisting}

\section{Interrupts}

The VGC IRQ helpers install the CPU IRQ vector, enable or disable source masks,
and acknowledge pending source bits. The source constants come from
\novacode{nova.inc}; the common masks are \novacode{VGC\_IRQ\_VBLANK},
\novacode{VGC\_IRQ\_COPPER0}, \novacode{VGC\_IRQ\_SPRCOLL}, and
\novacode{VGC\_IRQ\_SPRBG}.

\begin{lstlisting}
install_vblank_irq:
      LDA   #>vblank_irq
      LDY   #<vblank_irq
      JSR   vgc_irq_install

      LDA   #VGC_IRQ_VBLANK
      JSR   vgc_irq_enable
      RTS

vblank_irq:
      PHA
      LDA   VGC_IRQ_STATUS
      AND   #VGC_IRQ_VBLANK
      BEQ   @done

      ; Do one-frame work here.

      LDA   #VGC_IRQ_VBLANK
      JSR   vgc_irq_ack
@done:
      PLA
      RTI
\end{lstlisting}

Real handlers should preserve any registers they use and must acknowledge every
VGC source they service before returning. If a source remains pending and
enabled, the IRQ line remains asserted.

\section{Getting The Most From VGC}

Use \novacode{vgc\_vsync} as the frame boundary for animation. Update game
state, write sprite attributes, and start large drawing work in a consistent
order each frame.

Use direct sprite registers for frequent movement. The command helpers are
convenient, but \novacode{\$A040 + sprite * 8} is the faster path for X, Y,
shape, flags, priority, and transparent color.

Use the graphics primitives for small generated shapes and diagnostics. Use
DMA, blitter, NVG loading, or pager transfers for large images, sprite banks,
and full-screen changes.

Set \novacode{VGC\_GFXTRANS} deliberately. It defaults to color 0 so black
graphics pixels are transparent. If a program needs visible black pixels in the
graphics layer, move the transparent color to a palette index that the art does
not use and keep blitter color keys in sync with that choice.

Use the copper for scanline work instead of tight CPU timing loops. Copper
events are tied to the video beam and continue to run while the 6502 works on
the next frame.
